\uuid{VhO7}
\titre{Changement de variable/positivité}
\auteur{hanani ; mzari}
\organisation{exo7}
\licence{CC BY-NC-SA 4.0}
\datecreate{2020-01-20}
\niveau{L1}
\module{Analyse}
\chapitre{Calcul d'intégrales}
\sousChapitre{Autre}
\difficulte{2}
\theme{Calculs d'intégrales, changement de variable, positivité}
\exo7qcmid{lille-2.integrales.014}

\contenu{
\begin{question}
\qclassification{127.99}
\qtags{motcle=changement de variable/positivité}
Soit $\displaystyle f(x)=\frac{x}{x^2+1}$. Parmi les affirmations suivantes, cocher celles qui sont vraies :
\begin{answers}
\bad{$\displaystyle \int _0^2f(x)\, \mathrm{d}x=\int _0^{4}\frac{\mathrm{d}t}{t+1}$.}
\good{$\displaystyle \frac{1}{5}\int _0^{2}x\, \mathrm{d}x<\int _0^2f(x)\, \mathrm{d}x<\int _0^{2}x\, \mathrm{d}x$.}
\good{$\displaystyle \int _{-1}^1f(x)\, \mathrm{d}x=0$.}
\bad{$\displaystyle \int _0^2f(x)\, \mathrm{d}x=\ln 5$.}
\end{answers}
\begin{explanations}
Avec $t=x^2$, on a : $\mathrm{d}t=2x\, \mathrm{d}x$, $t(0)=0$, $t(2)=4$ et
$$\int _0^2f(x)\,\mathrm{d}x=\frac{1}{2}\int _0^{4}\frac{\mathrm{d}t}{t+1}=\frac{1}{2}\Big[\ln (t+1)\Big]^4_0=\frac{1}{2}\ln 5.$$
Pour tout $x\in ]0,2[$, on a : $\displaystyle \frac{x}{5}<f(x)<x$. Donc $\displaystyle \frac{1}{5}\int _0^{2}x\, \mathrm{d}x<\int _0^2f(x)\, \mathrm{d}x<\int _0^{2}x\, \mathrm{d}x$. Enfin, la fonction $f$ étant impaire sur $\Rr$, pour tout $a>0$, on a : $\displaystyle \int _{-a}^af(x)\, \mathrm{d}x=0$.
\end{explanations}
\end{question}
}
